
% 10. The Maximization of Fairness and the Convergence of Competition: A Single Symmetric Competitive Track.tex

\section{공정성의 극대화와 경쟁의 융합: 대칭적 단일 경쟁 트랙}
우리는 현대 사회를 인류 역사상 가장 진보한 시대로 찬미한다. 왕과 백성, 귀족과 평민 사이의 좁힐 수 없던 신분적 비대칭성은 혁명을 통해 무너졌고, 교육과 정치 참여의 기회는 보편화되었으며, 무엇보다 남성과 여성 사이의 사회적 역할 비대칭성이 지속적으로 감소해 왔다. 우리는 문명이 걸어온 이 장엄한 궤적을 '공정성의 확대'라 부르며, 차별이 철폐된 사회가 우리에게 더 평화롭고 조화로운 세상을 가져다줄 것이라 굳게 믿는다.

그러나 디스턴스(DISTANCE)의 서늘한 열역학적 렌즈를 들이대면, 우리가 '공정성'이라 믿어온 이 고결한 도덕적 성취의 진짜 물리적 실체가 폭로된다. 공정성은 결코 갈등을 종식시키는 평화의 깃발이 아니다. 공정성은 사회 내부에 존재하던 모든 장벽(비대칭성)들을 깎아내어, 서로 다른 영역에 격리되어 있던 모멘텀들을 완벽하게 동일한 무대 위로 쏟아붓는 가장 잔혹하고 숨 막히는 '평활화(Equalization)'의 과정이다.

이러한 물리학적 진실은 매우 뼈아픈 역설을 낳는다. 공정성은 결코 경쟁을 제거하지 않는다. 오히려 경쟁을 '보편화'한다. 과거의 계급 사회나 전통적 성별 분업 사회에서는 많은 사람들이 애초에 특정 경쟁에 참여할 기회조차 갖지 못했다. 신분 때문에, 혹은 성별 때문에 사회적 트랙은 철저히 분리되어 있었고, 사람들은 자신이 속한 작은 영역(로컬 아일랜드) 안에서만 제한적인 에너지를 소모하며 경쟁했다. 그러나 문명의 거대한 평활화 엔진이 신분, 제도, 성별이라는 모든 인위적 장벽(에너지 갭)들을 매끄럽게 깎아내 버리자, 격리되어 있던 수많은 사람들의 모멘텀이 단 하나의 거대한 믹서기 안으로 빨려 들어가기 시작한 것이다.

이 폭발적인 경쟁의 융합이 가장 선명하게 드러나는 지점이 바로 남성과 여성의 관계이다. 오랫동안 남성과 여성은 서로 다른 사회적 역할(비대칭성)을 수행해 왔다. 이것은 단순한 문화적 억압의 결과가 아니라, 혹독한 자연환경 속에서 근육의 힘과 생물학적 조건에 따라 생존의 모멘텀을 가장 효율적으로 해소하기 위해 빚어진 불가피한 기능적 분업이었다. 남성과 여성은 서로 다른 트랙 위에서, 서로 다른 방식으로 평가받으며 살아갔다.

그러나 산업혁명이 근육의 중요성을 무너뜨리고, 대중 교육이 지식의 접근성을 열었으며, 정보 기술이 물리적 장벽을 해체하면서 두 개의 트랙을 가로막고 있던 벽은 무참히 붕괴되었다. 나아가 인간의 육체 노동을 로봇이 대체하고, 반복적인 사고 노동마저 인공지능(AI)이 대신하게 되면서, 생물학적 성별이 경제적 생산성이나 사회적 가치와 직결되던 모든 물리적 근거들은 완벽하게 평활화되어 버렸다.

그 결과, 오늘날 남성과 여성은 완전히 융합된 '대칭적 단일 경쟁 트랙(Symmetric Single Competitive Track)' 위로 밀어 넣어졌다. 이제 남녀는 같은 학교에서 공부하고, 같은 시험을 치르며, 같은 직장에 지원하여 동일한 승진 체계 안에서 가장 맹렬하게 서로의 모멘텀을 부딪치며 경쟁한다. 과거처럼 서로의 다름을 인정하며 각자의 영역을 구축하던 시대는 끝났다. 현대 사회는 모든 비대칭성을 해체하고, 모두를 단 하나의 잣대로 평가하는 우주 역사상 가장 매끄럽고 평탄한 궁극의 대칭적 투기장을 완성해 낸 것이다.

많은 이들은 모든 차별의 벽이 허물어지고 남녀가 완벽히 동일한 선상에서 경쟁하게 된 이 상태를 문명의 위대한 승리라 부른다. 하지만 디스턴스의 렌즈는 바로 이 완벽한 평활화의 무대 위에서 발생하게 될 가장 거대하고 치명적인 물리학적 파열음을 경고한다.
사회적 잣대가 완벽하게 대칭화되고 모든 것이 공정해진 그 텅 빈 투기장 위에서, 인간의 그 어떤 정치 제도나 경제 정책으로도 깎아낼 수 없는 단 하나의 거대한 '생물학적 비대칭성'이 홀로 섬뜩하게 그 거대한 모습을 드러내기 시작한 것이다. 문명이 빚어낸 이 완벽한 대칭성의 무대 위에서, 오직 여성의 몸을 통해서만 이루어지는 '임신과 출산'이라는 선천적 비대칭성은 과연 어떻게 해석될 것인가?

과거의 사회는 교육, 정치, 직업, 경제 등 모든 곳에 불평등이라는 얼룩이 가득 묻어 있는 지저분한 유리창이었다. 그래서 특정한 비대칭성(출산의 부담) 하나가 특별히 부각되지 않았다. 그러나 문명의 맹렬한 평활화 엔진이 그 유리창을 완벽하게 투명하고 매끄럽게 닦아내 버리자, 그 텅 빈 무대 위에 홀로 남겨진 번식의 생물학적 비대칭성은 이제 견딜 수 없는 거대한 모순으로 폭발하기 시작한 것이다.

동일한 교실에서 공부하고, 동일한 잣대로 평가받으며, 동일한 트랙 위에서 맹렬하게 성공의 모멘텀을 다투도록 설계된 이 완벽한 대칭성의 투기장 안에서, 오직 여성만이 짊어져야 하는 임신과 출산이라는 생물학적 조건은 더 이상 고결한 희생이나 자연의 섭리가 아니다. 그것은 공정 경쟁이라는 문명의 가장 숭고한 규칙을 정면으로 위반하는 치명적인 '선천적 기회비용'으로 인식될 수밖에 없는 것이다.

이 지점에서 인류가 그토록 찬미해 온 '공정성'의 서늘한 우주적 역설이 장엄하게 완성된다. 저출산은 단순히 청년들이 가난하거나 이기적이기 때문에 발생하는 경제적 파업이 아니다. 그것은 인류가 수천 년간 피 흘려 닦아온 '완벽한 공정성(평활화)'이라는 문명의 맹렬한 궤적과, 결코 평탄해질 수 없는 인간의 '비대칭적 번식 구조'가 정면으로 충돌하며 뿜어내는 파괴적인 물리학적 파열음이다.

가장 고도로 발전하고 가장 공정해진 선진국일수록 아이의 울음소리가 멈추는 이유는 그들이 부유해서가 아니다. 그들이 우주 역사상 가장 맹렬하고 완벽하게 닦인 '대칭적 단일 경쟁 트랙'을 완성해 냈기 때문이다. 문명이 비대칭성을 매끄럽게 깎아내려 할수록, 번식이라는 엇갈린 디스턴스는 그 매끄러운 표면 위에서 산산조각 날 수밖에 없는 웅장한 비극적 결말을 맞이하고 있는 것이다.